\documentclass[11pt]{article}
\usepackage[a4paper,margin=25mm]{geometry}
\usepackage{amsmath,amssymb,amsthm,booktabs,microtype}
\usepackage[T1]{fontenc}
\usepackage{lmodern}
\usepackage[dvipsnames]{xcolor}
\usepackage[colorlinks=true,linkcolor=MidnightBlue,citecolor=MidnightBlue,urlcolor=MidnightBlue]{hyperref}
\usepackage{fancyhdr}
\pagestyle{fancy}\fancyhf{}
\fancyhead[L]{\small Singer-invariant blocking sets of PG(2,13)}
\fancyhead[R]{\small Dev D. Goyal}
\fancyfoot[C]{\thepage}
\setlength{\headheight}{14pt}
\hypersetup{pdftitle={Blocking sets of PG(2,13) invariant under a fixed-point-free Singer subgroup meet some line in five points},
pdfauthor={Dev D. Goyal},
pdfsubject={Fully autonomous agentic research},
pdfkeywords={blocking sets, projective planes, PG(2,13), Singer cycle, Erdos problem 1159, line intersections, Lean 4}}
\newtheorem{theorem}{Theorem}
\newtheorem{lemma}[theorem]{Lemma}
\newtheorem{proposition}[theorem]{Proposition}
\newtheorem{corollary}[theorem]{Corollary}
\theoremstyle{remark}
\newtheorem{remark}[theorem]{Remark}
\newcommand{\PG}{\mathrm{PG}(2,13)}
\newcommand{\F}{\mathbb{F}}
\newcommand{\Z}{\mathbb{Z}}
\title{\textbf{Blocking sets of $\mathbf{PG(2,13)}$ invariant under\\a fixed-point-free Singer subgroup\\meet some line in five points}}
\author{Dev D. Goyal}
\date{31 October 2026}
\begin{document}
\maketitle
\begin{abstract}
Let $H$ be the subgroup of order three of a Singer cycle group of the
projective plane $\PG$; it fixes no point. We show that every $H$-invariant
blocking set of $\PG$ meets some line in at least five points. The
$21$-point blocking set of Blokhuis, Brouwer and Wilbrink is invariant under
a conjugate of $H$ and meets every line in at most five points, so five is
the least possible largest line intersection of an $H$-invariant blocking
set. Every fixed-point-free collineation of order three of $\PG$ generates a
conjugate of $H$, so the statement holds for all of them. This settles one
symmetric case of the plane of order $13$ in Erd\H{o}s Problem~1159, which
asks whether every finite projective plane has a point set meeting each line
in at least one and at most $C$ points; a public computational campaign had
recorded this case as unresolved. The proof reduces the question to $61$
weighted line conditions on the orbits of $H$ and shows by exhaustive search
that they have no solution with intersections at most four. A Lean~4 project
proves the theorem for the specified subgroup, with $\PG$ modelled by
Mathlib's projectivization of $\F_{13}^3$; the search runs in the Lean kernel
and is covered by a soundness proof. An independent Python program repeats
the geometry, the witness checks and the search.
\end{abstract}

\section{Introduction and result}

A \emph{blocking set} of a projective plane is a set of points meeting every
line. Erd\H{o}s asked whether there is a constant $C$ such that every finite
projective plane has a point set $S$ with $1\le|S\cap\ell|\le C$ for all lines
$\ell$ \cite{EP1159}; Erd\H{o}s, Silverman and Stein proved that
$|S\cap\ell|=O(\log n)$ can be achieved in a plane of order $n$
\cite{ESS}. For a particular plane the question becomes the computation of
\[
  c(\Pi)=\min_{B}\ \max_{\ell}\ |B\cap\ell|,
\]
the minimum over blocking sets $B$ of the largest line intersection. For the
Desarguesian plane $\PG$ of order $13$ the $21$-point blocking set of
Blokhuis, Brouwer and Wilbrink \cite{BBW} shows $c(\PG)\le5$. A public
computational campaign on this problem \cite{campaign} searched for blocking
sets of $\PG$ with all line intersections at most four that are invariant
under various collineation groups; for the subgroup of order three of a
Singer group, whose $61$ point orbits all have size three, its solver run
ended with status UNKNOWN after 90 seconds. We settle that case.

Let $\F=\F_{13}[\theta]/(\theta^3+\theta+6)$. The powers of $\theta$ are
pairwise distinct up to $\theta^{2196}=1$, so $\F$ is the field with $13^3$
elements and $\theta$ is primitive. We identify the points of $\PG$ with the
one-dimensional $\F_{13}$-subspaces of $\F$, written in the basis
$1,\theta,\theta^2$, and the lines with the two-dimensional subspaces.
Multiplication by $\theta$ induces a collineation $\sigma$ of order
$183=13^2+13+1$, a \emph{Singer cycle}; let $H=\langle\sigma^{61}\rangle$, of
order three. Since $\sigma^{61}$ maps the point $\langle\theta^i\rangle$ to
$\langle\theta^{i+61}\rangle\ne\langle\theta^i\rangle$, $H$ fixes no point.

\begin{theorem}\label{thm:main}
Every $H$-invariant blocking set of $\PG$ meets some line in at least five
points. There is an $H$-invariant minimal blocking set of $21$ points that
meets every line in at most five points. Hence
\[
  \min\bigl\{\max_{\ell}|B\cap\ell| : B \text{ an $H$-invariant blocking set of } \PG\bigr\}=5 .
\]
\end{theorem}

The lower bound has no size or minimality assumption on $B$. The upper bound
is the known example of \cite{BBW} (Section~\ref{sec:witness}). Since
$13$ is prime, every collineation of $\PG$ is a projectivity.

\begin{corollary}\label{cor:all}
Let $g$ be a collineation of $\PG$ of order three without fixed points. Then
every $\langle g\rangle$-invariant blocking set meets some line in at least
five points, and five is attained.
\end{corollary}

The corollary follows from Theorem~\ref{thm:main} and the conjugacy argument
of Section~\ref{sec:conj}. It does not decide whether $c(\PG)$ is $4$ or $5$:
blocking sets without this symmetry are not covered.

\section{The cyclic model and the reduction}

For $i\in\Z_{183}$ let $P_i=\langle\theta^i\rangle$; these are all $183$
points. The points whose coordinate vector has constant coefficient zero form
a line,
\[
  L_0=\{P_d : d\in D\},\qquad D=\{1,2,4,17,24,29,43,77,83,87,120,138,155,176\},
\]
and the lines of $\PG$ are exactly $L_t=\sigma^t(L_0)=\{P_{d+t}:d\in D\}$ for
$t\in\Z_{183}$; $D$ is a planar difference set modulo $183$. The subgroup $H$
acts by $i\mapsto i+61$, with the $61$ orbits $O_j=\{P_j,P_{j+61},P_{j+122}\}$,
$j\in\Z_{61}$. An $H$-invariant set is $B_S=\bigcup_{j\in S}O_j$ for a unique
$S\subseteq\Z_{61}$, and
\[
  |B_S\cap L_t| = r_S(t) := \#\{d\in D : d+t \bmod 61\in S\},
\]
which depends only on $t\bmod 61$. Modulo $61$ the elements of $D$ are
\[
  1,2,4,16,16,17,22,24,26,29,33,43,54,59,
\]
with $16$ twice ($77$ and $138$) and the others once. Thus each of the $61$
conditions involves twelve orbits with weight one and one orbit, the
\emph{doubled} orbit $t+16$, with weight two. The $61\times61$ weight matrix
$M$ is circulant and satisfies $MM^{\mathsf T}=13I+3J$.

\begin{lemma}\label{lem:reduction}
For $c\ge1$ there is an $H$-invariant blocking set with all line intersections
at most $c$ if and only if there is $S\subseteq\Z_{61}$ with $0\in S$ and
$1\le r_S(t)\le c$ for all $t\in\Z_{61}$.
\end{lemma}

\begin{proof}
The first condition holds for $B_S$ exactly when $1\le r_S(t)\le c$ for all
$t$. A blocking set is nonempty, so some $j\in S$. The collineation
$\sigma^{-j}$ commutes with $H$, maps $B_S$ to $B_{S-j}$ and permutes the
lines, so $S-j$ is again a solution, and it contains $0$.
\end{proof}

\begin{remark}
Since $\sum_t r_S(t)=14|S|$ and, by $M^{\mathsf T}M=13I+3J$,
$\sum_t r_S(t)^2=3|S|^2+13|S|$, summing $(r_S(t)-1)(r_S(t)-4)\le0$ shows that a
solution for $c=4$ would have $7\le|S|\le12$. This is not used below.
\end{remark}

\section{No solution with intersections at most four}\label{sec:search}

\begin{proposition}\label{prop:unsat}
There is no $S\subseteq\Z_{61}$ with $0\in S$ and $1\le r_S(t)\le4$ for all
$t\in\Z_{61}$.
\end{proposition}

\begin{proof}
We search over partial assignments, which fix some orbits as chosen and some
as excluded; orbit $0$ is chosen at the start. A node is closed when some
condition $t$ already has weight at least five on the chosen orbits, or all
its orbits are excluded. Before branching we apply the forced steps: if a
condition with no chosen orbit has exactly one orbit that is not excluded,
that orbit is chosen; if a condition has weight four, its other orbits are
excluded; if it has weight three, its doubled orbit is excluded. Each step
is valid for every solution extending the node. If no node is closed and no
step applies, we pick a condition with no chosen orbit and fewest available
orbits and branch on which of them is the first chosen one. Every branch of
this finite search closes, so no solution exists. The Python implementation
described in Section~\ref{sec:verification} visits $156{,}717$ nodes.
\end{proof}

The lower bound of Theorem~\ref{thm:main} follows from Lemma~\ref{lem:reduction}
with $c=4$ and Proposition~\ref{prop:unsat}.

\section{The matching example}\label{sec:witness}

Let $S_0=\{0,1,13,14,47,48,60\}\subseteq\Z_{61}$. The $21$-point set
$B_{S_0}$ meets $126$ lines in one point, $36$ lines in two points, $9$ lines
in four points and $12$ lines in five points; every point of $B_{S_0}$ lies on
a tangent line, so $B_{S_0}$ is a minimal blocking set. This proves the upper
bound in Theorem~\ref{thm:main}.

The set is the sporadic example of \cite[Section~5.3]{BBW}. In the
coordinates of that paper it consists of the three vertices of the reference
triangle, the points $(1,a,0)$, $(0,1,a)$, $(a,0,1)$ with $a^3=-1$, and the
points $(b,c,1)$ with $b^3=c^3=1$. It is invariant under the projectivity of
\[
  A=\begin{pmatrix}0&12&0\\0&0&4\\12&0&0\end{pmatrix},\qquad A^3=4I,
\]
which has no fixed point because $4$ is not a cube in $\F_{13}$. With
\[
  K=\begin{pmatrix}1&7&8\\0&10&1\\0&4&12\end{pmatrix}
\]
one has $K(7A)=TK$, where $T$ is the matrix of multiplication by
$\theta^{61}$, and $K$ maps the set of \cite{BBW} onto $B_{S_0}$. No new
construction is claimed.

\section{All fixed-point-free subgroups of order three}\label{sec:conj}

Let $g$ be a collineation of order three without fixed points and let $A$ be
a matrix inducing it, so $A^3=aI$ with $a\in\F_{13}^\times$. If $a=b^3$, then
$(b^{-1}A)^3=I$; since $13\equiv1\pmod3$, the polynomial $X^3-1$ splits into
distinct linear factors over $\F_{13}$, so $b^{-1}A$ is diagonalizable and
has an eigenvector, a fixed point of $g$. Hence $a$ is not a cube, $X^3-a$ has
no root and is irreducible, and it is the minimal polynomial of the
nonscalar matrix $A$. Then $E=\F_{13}[A]$ is a field with $13^3$ elements and
$\F_{13}^3$ is a one-dimensional $E$-vector space, so $g$ lies in the image of
$E^\times$ in $\mathrm{PGL}(3,13)$, a Singer group. Any two subfields of the
matrix algebra isomorphic to $\F_{13^3}$ are conjugate under $\mathrm{GL}(3,13)$
(identify $\F_{13}^3$ with either field), so all Singer groups are conjugate,
and a cyclic group of order $183$ has exactly one subgroup of order three.
Therefore $\langle g\rangle$ is conjugate to $H$. A projectivity carrying $H$
to $\langle g\rangle$ carries $H$-invariant blocking sets to
$\langle g\rangle$-invariant ones and preserves line intersections, which
gives Corollary~\ref{cor:all}. This section is a conventional argument and is
not part of the Lean proof.

\section{Verification}\label{sec:verification}

\paragraph{Lean.} The package contains a Lean~4.35.0 project using Mathlib
revision \texttt{56a0058}. Points and lines of $\PG$ are both elements of
Mathlib's \texttt{Projectivization (ZMod 13) (Fin 3 \(\to\) ZMod 13)}, and a
point lies on a line when the dot product of representatives vanishes. The
Singer cycle is the projectivity of the matrix of multiplication by $\theta$,
and $H$ is generated by its $61$st iterate. The project proves that the cycle
has order exactly $183$, that its $61$st iterate has order three and fixes no
point, the properties of the $21$-point witness, the reduction of
Lemma~\ref{lem:reduction} for arbitrary point sets, Proposition~\ref{prop:unsat},
and the final statement \texttt{SingerC3.isLeast\_admissibleCaps}: $5$ is the
least $c$ for which an $H$-invariant blocking set with all line intersections
at most $c$ exists.

The search of Proposition~\ref{prop:unsat} is a Lean function \texttt{go}
that is evaluated by the kernel with \texttt{decide +kernel}. It takes only
five kinds of steps, each checked against the actual line counts $r_S$:
closing a node, excluding an orbit after checking that choosing it would close
the node, choosing an orbit after checking that excluding it would close the
node, excluding the available orbits of a condition of weight four, and
splitting on an orbit. A separate heuristic on packed line counts decides
which step to take. The soundness theorem \texttt{go\_sound} holds for every
value of the heuristic, so its correctness is not needed. The search is split
into a top search from the root and $342$ subsearches with at most $4{,}961$
steps each; together they take $771{,}752$ steps. Finite facts about the
coordinate model are also decided by the kernel. There is no
\texttt{native\_decide}, no compiled-code evaluation and no added axiom, and
every declaration of the project depends only on \texttt{propext},
\texttt{Classical.choice} and \texttt{Quot.sound}. The kernel search takes
about an hour of processor time in total.

\paragraph{Python.} The program \texttt{certificate/check\_singer.py} uses only
the Python standard library. It recomputes the field, the $183$ points, all
lines from their coordinate equations, the difference set $D$, the $61$
conditions and the identity $MM^{\mathsf T}=13I+3J$; it checks the witness,
its intersection numbers, its tangents and the conjugacy to the set of
\cite{BBW}; and it runs the search of Proposition~\ref{prop:unsat}.

\section{Limitations}

The result concerns one conjugacy class of subgroups. Whether $\PG$ has a
blocking set with all line intersections at most four, without symmetry, is
not decided here. The identification of $D$ as the line set is checked by
computation and in Lean; the conjugacy statement of Section~\ref{sec:conj} is
not formalized. The literature comparison is bounded, and historical priority
is not certified.

\section*{Research process and AI disclosure}

This is fully autonomous agentic research. Dev D. Goyal is responsible for
publication. No expert validation or external peer review is claimed, and
the novelty assessment is provisional.

Autonomous AI research agents selected the question, found the reduction to
$61$ orbit conditions, identified the witness with the example of
\cite{BBW}, produced and checked a solver certificate for
Proposition~\ref{prop:unsat}, wrote their own Lean proof and searched the
literature in an earlier session, which also ran internal AI-assisted reviews
that are not peer review. Its records describe the research and review
agents as Astra agents. The same session proposed a second result, a
logarithmic lower bound for the number of bounded-degree plane curves whose
rational points block all lines; it is not part of this paper.

Opus 5.5 by Anthropic used with Claude Code assessed the package for this
release. It wrote the Lean project used here, including the kernel search and
its soundness proof, ran its fresh build and axiom audit, wrote the
independent Python checker and this manuscript, and integrated the package
into the repository. The solver certificate and the earlier Lean proof are
not included. AI checks are not human peer review.

\begin{thebibliography}{9}
\bibitem{BBW}
A.~Blokhuis, A.~E.~Brouwer and H.~A.~Wilbrink, \emph{Blocking sets in
PG$(2,p)$ for small $p$, and partial spreads in PG$(3,7)$}, Advances in
Geometry \textbf{3} (2003), special issue, S245--S253.
\href{https://doi.org/10.1515/advg.2003.2003.s1.245}{doi:10.1515/advg.2003.2003.s1.245}.

\bibitem{EP1159}
T.~F.~Bloom, \emph{Erd\H{o}s Problem \#1159},
\href{https://www.erdosproblems.com/1159}{erdosproblems.com/1159},
accessed 31 October 2026.

\bibitem{ESS}
P.~Erd\H{o}s, R.~Silverman and A.~Stein, \emph{Intersection properties of
families containing sets of nearly the same size}, Ars Combinatoria
\textbf{15} (1983), 247--259.

\bibitem{campaign}
Erd\H{o}s problem verification campaign, problem 1159, orbit-search log for
$q=13$, commit \texttt{7bccace},
\href{https://github.com/pw/erdosproblemaday-verification/blob/7bccace07c6b555623be274558e463b5ef2c3ea1/1159/1159-klein-blocks-pg2-11-and-the-q13-question/orbitzoo/orbitzoo_q13.log}{github.com/pw/erdosproblemaday-verification},
accessed 31 October 2026.
\end{thebibliography}
\end{document}
